\documentclass{article}
\usepackage[T1]{fontenc}
\usepackage{lmodern}
\usepackage{graphicx}
\usepackage{amsmath,amssymb}
\usepackage[a4paper,margin=2cm]{geometry}
\usepackage[absolute,overlay]{textpos}
\setlength{\TPHorizModule}{1mm}
\setlength{\TPVertModule}{1mm}
\usepackage[indonesian]{babel}
\usepackage{hyperref}
\setlength{\emergencystretch}{2em}
% unit: Halaman 14

\begin{document}

% === Halaman 14 (llm) ===

\begin{itemize}
    \item Prosedur membuat kunci publik dan kunci privat:
\end{itemize}

\begin{enumerate}
    \item Tentukan barisan \textit{superincreasing}.
    \item Kalikan setiap elemen di dalam barisan tersebut dengan $n \pmod{m}$ \\ ( Modulus $m$ seharusnya angka yang lebih besar daripada jumlah semua elemen di dalam barisan, sedangkan pengali $n$ seharusnya tidak mempunyai faktor persekutuan dengan $m$, atau $\text{PBB}(n, m) = 1$ )
    \item Hasil perkalian akan menjadi kunci publik sedangkan barisan \textit{superincreasing} semula menjadi kunci privat.
\end{enumerate}

\end{document}
